% !TEX root = ../main.tex
\begin{ouverture}[gazparfait]{18}{Modèle du gaz parfait}{Partie 3 · L'énergie : conversions et transferts}
{Échelles microscopique et macroscopique\enspace·\enspace Équation d'état\enspace·\enspace Zéro absolu\enspace·\enspace Limites du modèle}
\objectif{Relier la masse volumique, la pression et la température d'un gaz aux propriétés de ses entités à l'échelle
microscopique.}
\objectif{Exploiter l'équation d'état du gaz parfait pour décrire le comportement d'un gaz.}
\objectif{Calculer la quantité de matière, le volume molaire et la masse volumique d'un gaz.}
\objectif{Identifier quelques limites du modèle du gaz parfait.}
\objectif{Tester l'équation d'état à partir de mesures, à l'aide d'un programme Python.}
\prerequis{Pression, force pressante, loi de Mariotte (classe de première).}
\prerequis{Quantité de matière, masse molaire, constante d'Avogadro (classe de seconde).}
\prerequis{Statique des fluides et poussée d'Archimède (chapitre~17).}
\end{ouverture}
\essentiel
\essSection{1}{Du microscopique au macroscopique}
\begin{definition}[Décrire un gaz]
Un gaz est constitué d'un très grand nombre d'entités (molécules ou atomes) en mouvement incessant et désordonné :
c'est l'\textbf{agitation thermique}. À l'échelle macroscopique, son état est décrit par des grandeurs mesurables :
le volume $V$, la masse volumique $\rho=m/V$, la pression $p$ et la \textbf{température thermodynamique} $T$, exprimée en
kelvins : $T=\theta+273{,}15$, avec $\theta$ en degrés Celsius.
\end{definition}
\begin{propriete}[Interprétation microscopique]
\begin{itemize}
\item La masse volumique dépend du nombre d'entités par unité de volume et de leur masse.
\item La pression résulte des chocs des entités sur les parois : elle augmente avec le nombre de chocs par seconde et
avec leur violence.
\item La température mesure l'agitation des entités : elle augmente avec leur énergie cinétique moyenne. Au
\textbf{zéro absolu}, $T=\qty{0}{K}$, soit \qty{-273,15}{\degreeCelsius}, l'agitation serait minimale ; aucune
température ne lui est inférieure.
\end{itemize}
\end{propriete}
\schema{18-agitation}
\begin{exemple}
Un litre d'air à \qty{20}{\degreeCelsius} sous \qty{1,00}{bar} contient \qty{0,041}{mol} de molécules, soit environ
\num{2.5e22} molécules, animées de vitesses de l'ordre de \qty{500}{m.s^{-1}}. Chaque centimètre carré de paroi reçoit
des milliards de milliards de chocs par seconde : à notre échelle, leur effet moyen est une force pressante constante.
\end{exemple}
\essSection{2}{L'équation d'état du gaz parfait}
\begin{definition}[Modèle du gaz parfait]
Dans le modèle du \textbf{gaz parfait}, les entités sont assimilées à des points matériels, sans interaction entre
elles en dehors des chocs. Ce modèle décrit bien les gaz réels peu denses.
\end{definition}
\begin{loi}[Équation d'état du gaz parfait]
Pour une quantité de matière $n$ de gaz parfait, de pression $p$, de volume $V$ et de température $T$ :
\[p\,V=n\,R\,T,\]
avec $p$ en pascals, $V$ en \unit{m^3}, $n$ en moles, $T$ en kelvins et $R=\qty{8,314}{J.mol^{-1}.K^{-1}}$, la constante des gaz
parfaits.
\end{loi}
\begin{propriete}[Conséquences]
\begin{itemize}
\item À $n$ et $T$ fixés, $p\,V$ est constant : c'est la loi de Mariotte, vue en première.
\item À $n$ et $V$ fixés, $p$ est proportionnelle à $T$ ; à $n$ et $p$ fixés, $V$ est proportionnel à $T$.
\item Le volume molaire $V_\mathrm{m}=\dfrac{R\,T}{p}$ est le même pour tous les gaz parfaits : \qty{22,4}{L.mol^{-1}} à
\qty{0}{\degreeCelsius} et \qty{24,1}{L.mol^{-1}} à \qty{20}{\degreeCelsius}, sous \qty{1013}{hPa}.
\item La masse volumique d'un gaz de masse molaire $M$ vaut $\rho=\dfrac{m}{V}=\dfrac{n\,M}{V}=\dfrac{p\,M}{R\,T}$.
\end{itemize}
\end{propriete}
\begin{savoirfaire}[Exploiter l'équation d'état]
Préciser le système et vérifier si sa quantité de matière reste constante. Convertir en unités SI :
$\qty{1}{L}=\qty{e-3}{m^3}$, $\qty{1}{bar}=\qty{e5}{Pa}$, $\qty{1}{hPa}=\qty{100}{Pa}$, $T=\theta+273{,}15$. Si $n$ ne change pas
entre deux états 1 et 2 :
\[\frac{p_1\,V_1}{T_1}=\frac{p_2\,V_2}{T_2}\qquad(=n\,R).\]
\end{savoirfaire}
\begin{exemple}
L'air ($M=\qty{29,0}{g.mol^{-1}}$) à \qty{20}{\degreeCelsius} sous \qty{1013}{hPa} a pour masse volumique
$\rho=\dfrac{\num{1.013e5}\times\num{0.0290}}{8{,}314\times293{,}15}=\qty{1,21}{kg.m^{-3}}$. Chauffé à volume constant, il
garde la même masse volumique, mais sa pression augmente.
\end{exemple}
\schema{18-zero}
\essSection{3}{Les limites du modèle}
\begin{propriete}[Domaine de validité]
Le modèle du gaz parfait est valable lorsque le gaz est \textbf{peu dense} : pressions modérées, de l'ordre de
quelques bars au plus, et températures éloignées de celle de liquéfaction. À haute pression, le volume propre des
molécules n'est plus négligeable ; à basse température, leurs interactions attractives deviennent importantes et le
gaz finit par se liquéfier, ce que le modèle ne prévoit pas.
\end{propriete}
\begin{exemple}
À \qty{0}{\degreeCelsius} sous \qty{1013,25}{hPa}, le volume molaire du dioxyde de carbone vaut \qty{22,26}{L.mol^{-1}}, au lieu de
\qty{22,41}{L.mol^{-1}} pour un gaz parfait : l'écart n'est que de \qty{0,7}{\percent}. Mais à \qty{20}{\degreeCelsius}, le
dioxyde de carbone se liquéfie dès \qty{57}{bar} : dans un extincteur, il est liquide.
\end{exemple}
\begin{codecote}[0.6]
\begin{pythoncol}
V = [60.0, 50.0, 40.0, 30.0, 20.0]   # mL
p = [1013, 1218, 1521, 2024, 3035]   # hPa
R, T = 8.314, 293.0                  # SI, K
for k in range(5):
    n = p[k]*1e2 * V[k]*1e-6 / (R*T)
    print(V[k], round(p[k] * V[k]),
          round(1000 * n, 3))
\end{pythoncol}
\begin{consolecol}
60.0 60780 2.495
50.0 60900 2.5
40.0 60840 2.498
30.0 60720 2.493
20.0 60700 2.492
\end{consolecol}
\cote
\fontsize{9.5}{12.5}\selectfont
\capacite{Capacité numérique}\par\smallskip
Les listes contiennent les volumes d'air d'une seringue et les pressions mesurées par un capteur, à
\qty{20}{\degreeCelsius}. Le programme affiche le produit $p\,V$, en \unit{hPa.mL}, et la quantité de matière
$n=\dfrac{p\,V}{R\,T}$, en millimoles : les deux restent constants, aux incertitudes de mesure près. L'air de la seringue
se comporte comme un gaz parfait.
\end{codecote}
\begin{attention}
L'équation d'état s'utilise avec la température en \textbf{kelvins} et la pression \textbf{absolue}. Le manomètre d'une
station-service indique une pression relative, c'est-à-dire l'écart avec la pression atmosphérique : un pneu « gonflé à
\qty{2,0}{bar} » contient de l'air à environ \qty{3,0}{bar}.
\end{attention}
\begin{experience}[Tester la loi de Mariotte et l'équation d'état]
On relie une seringue à un capteur de pression piloté par un microcontrôleur. À température constante, on fait varier
le volume d'air enfermé : la pression est inversement proportionnelle au volume, $p$ en fonction de $1/V$ donne une
droite passant par l'origine, de coefficient directeur $n\,R\,T$. On plonge ensuite un ballon fermé dans des bains d'eau
de températures croissantes : à volume constant, la pression varie linéairement avec la température, et la droite
prolongée s'annule vers \qty{-273}{\degreeCelsius}.
\end{experience}
\begin{rediger}
Préciser le système (le gaz, dont la quantité de matière est fixée ou non) et écrire l'équation d'état pour chaque état.
Convertir toutes les grandeurs en unités SI \textbf{avant} le calcul, puis contrôler l'ordre de grandeur : une mole de
gaz occupe environ \qty{24}{L} dans les conditions usuelles.
\end{rediger}
\exercices

\rubrique{Découverte}
\begin{blocLivre}{ex-18-1}
\begin{exercice}[1]{Interpréter à l'échelle microscopique}
Un gaz est enfermé dans une seringue bouchée.
\begin{questions}
\item Quelles grandeurs macroscopiques permettent de décrire l'état du gaz ? Pour chacune, donner une interprétation
à l'échelle des molécules.
\item On enfonce le piston à température constante. Comment évoluent la masse volumique et la pression du gaz ?
Interpréter à l'échelle microscopique.
\item On chauffe la seringue en bloquant le piston. Comment évoluent la pression et la masse volumique ? Interpréter.
\item Convertir en kelvins les températures \qty{20}{\degreeCelsius}, \qty{-18}{\degreeCelsius} et \qty{100}{\degreeCelsius}. Pourquoi une
température de \qty{-300}{\degreeCelsius} n'a-t-elle aucun sens ?
\end{questions}
\end{exercice}
\end{blocLivre}
\begin{blocLivre}{ex-18-2}
\begin{exercice}[1]{L'air d'une chambre}
Une chambre mesure \qty{4,0}{m} sur \qty{5,0}{m}, pour une hauteur de \qty{2,5}{m}. L'air y est à \qty{20}{\degreeCelsius}
sous \qty{1013}{hPa} ; on l'assimile à un gaz parfait de masse molaire \qty{29,0}{g.mol^{-1}}.
\begin{donnees}
$R=\qty{8,314}{J.mol^{-1}.K^{-1}}$.
\end{donnees}
\begin{questions}
\item Calculer le volume molaire d'un gaz parfait à \qty{0}{\degreeCelsius}, puis à \qty{20}{\degreeCelsius}, sous \qty{1013}{hPa}.
\item Calculer la quantité de matière d'air contenue dans la chambre.
\item En déduire la masse de cet air. Le résultat vous surprend-il ?
\item Le chauffage porte l'air à \qty{25}{\degreeCelsius}, la pression restant égale à \qty{1013}{hPa}. Quelle quantité d'air
quitte la pièce ? Quelle est sa masse ?
\end{questions}
\end{exercice}
\end{blocLivre}
\begin{blocLivre}{ex-18-3}
\begin{exercice}[1]{La seringue bouchée}
Une seringue bouchée contient \qty{50}{mL} d'air à la pression atmosphérique, \qty{1013}{hPa}. La section de son piston
vaut \qty{2,5}{cm^2}. La température reste constante.
\begin{questions}
\item Quelle loi relie la pression et le volume de l'air enfermé ? Justifier son utilisation.
\item On enfonce le piston jusqu'à la graduation \qty{20}{mL}. Calculer la nouvelle pression.
\item Quelle force faut-il alors exercer sur le piston, l'air extérieur exerçant la pression atmosphérique sur son
autre face ?
\item On tire maintenant le piston jusqu'à \qty{80}{mL}. Calculer la pression de l'air enfermé.
\end{questions}
\end{exercice}
\end{blocLivre}
\begin{blocLivre}{ex-18-4}
\begin{exercice}[1]{Le pneu qui chauffe}
Avant un départ en vacances, un conducteur règle la pression de ses pneus à \qty{2,30}{bar}, valeur lue sur le manomètre
de la station, par une température de \qty{10}{\degreeCelsius}. Après deux heures d'autoroute, l'air des pneus atteint
\qty{50}{\degreeCelsius}. On prend $p_\mathrm{atm}=\qty{1,00}{bar}$ et l'on néglige la variation du volume des pneus.
\begin{questions}
\item Le manomètre indique l'écart entre la pression de l'air du pneu et la pression atmosphérique. Quelle est la
pression absolue de l'air dans le pneu au départ ?
\item Calculer la pression absolue, puis la pression lue sur le manomètre, après deux heures de route.
\item De quel pourcentage la pression absolue a-t-elle augmenté ?
\item Pourquoi recommande-t-on de vérifier la pression des pneus « à froid » ?
\end{questions}
\end{exercice}
\end{blocLivre}
\begin{blocLivre}{ex-18-5}
\begin{exercice}[1]{La masse volumique d'un gaz}
On considère de l'air ($M=\qty{29,0}{g.mol^{-1}}$), de l'hélium ($M=\qty{4,0}{g.mol^{-1}}$) et du dioxyde de carbone
($M=\qty{44,0}{g.mol^{-1}}$), à \qty{20}{\degreeCelsius} sous \qty{1013}{hPa}.
\begin{donnees}
$R=\qty{8,314}{J.mol^{-1}.K^{-1}}$.
\end{donnees}
\begin{questions}
\item Établir l'expression $\rho=\dfrac{p\,M}{R\,T}$ de la masse volumique d'un gaz parfait.
\item Calculer la masse volumique de chacun des trois gaz.
\item Pourquoi un ballon gonflé à l'hélium s'élève-t-il dans l'air ?
\item Lors de la fermentation du raisin, du dioxyde de carbone se dégage. Pourquoi est-il dangereux de descendre dans
une cuve de vinification ?
\end{questions}
\end{exercice}
\end{blocLivre}
\rubrique{Entraînement}
\begin{blocLivre}{ex-18-6}
\begin{exercice}[2]{La bouteille de plongée}
Une bouteille de plongée de \qty{12}{L} est remplie d'air sous \qty{200}{bar} (pression absolue), à \qty{20}{\degreeCelsius}.
Le plongeur descend à \qty{20}{m}, où l'eau est à \qty{12}{\degreeCelsius} et la pression vaut \qty{3,0}{bar}. Il respire
\qty{20}{L} d'air par minute, volume mesuré à la pression ambiante.
\begin{donnees}
$R=\qty{8,314}{J.mol^{-1}.K^{-1}}$ ; on assimile l'air à un gaz parfait.
\end{donnees}
\begin{questions}
\item Calculer la quantité d'air contenue dans la bouteille.
\item Quel volume cet air occuperait-il sous \qty{1}{bar}, à \qty{20}{\degreeCelsius} ?
\item La bouteille se met à la température de l'eau. Calculer la nouvelle pression de l'air qu'elle contient.
\item Le plongeur doit remonter quand la pression dans la bouteille atteint \qty{50}{bar}. Calculer son autonomie à
\qty{20}{m} de profondeur.
\end{questions}
\end{exercice}
\end{blocLivre}
\begin{blocLivre}{ex-18-7}
\begin{exercice}[2]{La bulle du plongeur}
À \qty{20}{m} de profondeur, un plongeur expire une bulle d'air de \qty{1,0}{cm^3}. L'eau de mer est à \qty{12}{\degreeCelsius}
au fond et à \qty{18}{\degreeCelsius} en surface.
\begin{donnees}
$p_\mathrm{atm}=\qty{1,013e5}{Pa}$ ; $\rho_\mathrm{mer}=\qty{1025}{kg.m^{-3}}$ ; $g=\qty{9,81}{N.kg^{-1}}$ ; loi de la statique des
fluides : $p=p_\mathrm{atm}+\rho\,g\,h$ à la profondeur $h$.
\end{donnees}
\begin{questions}
\item Calculer la pression de l'eau à \qty{20}{m} de profondeur.
\item En supposant que la bulle garde la même quantité d'air, calculer son volume lorsqu'elle atteint la surface.
\item Quelle part de cette augmentation est due au changement de température ?
\item Un plongeur en bouteille a les poumons remplis de \qty{6,0}{L} d'air à \qty{20}{m}. Que deviendrait ce volume s'il
remontait en bloquant sa respiration ? En déduire une règle de sécurité.
\end{questions}
\end{exercice}
\end{blocLivre}
\begin{blocLivre}{ex-18-8}
\begin{exercice}[2]{Tester l'équation d'état en Python}
On reprend le programme de l'essentiel, qui exploite des mesures réalisées avec une seringue et un capteur de pression.
\begin{questions}
\item \capacite{Capacité numérique}Justifier les facteurs \code{1e2} et \code{1e-6} dans le calcul de \code{n}. Que
représentent les deux nombres calculés pour chaque mesure ?
\item Pourquoi ces résultats valident-ils la loi de Mariotte ? Les écarts observés sont-ils significatifs ?
\item Écrire les lignes qui tracent $p$ en fonction de $1/V$. Quelle courbe obtient-on ? Que vaut, en joules, son
coefficient directeur ?
\item On réchauffe la seringue à \qty{35}{\degreeCelsius}, le volume étant bloqué à \qty{20}{mL}. Quelle pression le capteur
doit-il indiquer ?
\end{questions}
\end{exercice}
\end{blocLivre}
\begin{blocLivre}{ex-18-9}
\begin{exercice}[2]{À la recherche du zéro absolu}
Un ballon de verre fermé, de volume constant, contient de l'air relié à un capteur de pression. On le plonge dans des
bains de températures différentes.
\begin{center}
\begin{tabular}{lccccc}
\toprule
$\theta$ (\unit{\degreeCelsius}) & 10 & 30 & 50 & 70 & 90 \\
$p$ (hPa) & 881 & 944 & \num{1005} & \num{1068} & \num{1130} \\
\bottomrule
\end{tabular}
\end{center}
\begin{questions}
\item Tracer l'allure de $p$ en fonction de $\theta$. Les mesures sont-elles compatibles avec une fonction affine ?
\item À l'aide des mesures extrêmes, déterminer le coefficient directeur et l'ordonnée à l'origine de la droite.
\item Pour quelle température la pression s'annulerait-elle ? Comparer à la valeur attendue.
\item Montrer que, si l'on utilise la température thermodynamique $T$, la pression est proportionnelle à $T$, comme le
prévoit l'équation d'état.
\end{questions}
\end{exercice}
\end{blocLivre}
\begin{blocLivre}{ex-18-10}
\begin{exercice}[2]{La bouteille écrasée}
Une bouteille en plastique de \qty{1,5}{L}, vide, est fermée à \qty{25}{\degreeCelsius} sous \qty{1013}{hPa}, puis placée dans un
congélateur à \qty{-18}{\degreeCelsius}.
\begin{questions}
\item En supposant le volume de la bouteille constant, calculer la pression finale de l'air qu'elle contient.
\item Calculer l'écart entre la pression extérieure et la pression intérieure.
\item Une face de la bouteille a une surface de \qty{1,0e-2}{m^2}. Calculer la force pressante résultante qu'elle subit et
la comparer au poids d'une masse que l'on donnera.
\item La bouteille s'écrase. Quel volume l'air occuperait-il si sa pression restait égale à \qty{1013}{hPa} ?
\end{questions}
\end{exercice}
\end{blocLivre}
\begin{blocLivre}{ex-18-11}
\begin{exercice}[2]{Les limites du modèle}
Le tableau donne les volumes molaires de quelques gaz mesurés à \qty{0}{\degreeCelsius} sous \qty{1013,25}{hPa}.
\begin{center}
\begin{tabular}{lcccc}
\toprule
Gaz & \ce{N2} & \ce{O2} & \ce{CO2} & \ce{NH3} \\
$V_\mathrm{m}$ (\unit{L.mol^{-1}}) & 22,40 & 22,39 & 22,26 & 22,09 \\
\bottomrule
\end{tabular}
\end{center}
\begin{donnees}
$R=\qty{8,314}{J.mol^{-1}.K^{-1}}$ ; $M(\ce{CO2})=\qty{44,0}{g.mol^{-1}}$ ; à \qty{20}{\degreeCelsius}, le dioxyde de carbone se liquéfie
sous \qty{57}{bar}.
\end{donnees}
\begin{questions}
\item Calculer le volume molaire d'un gaz parfait dans ces conditions, avec quatre chiffres significatifs.
\item Calculer l'écart relatif pour chaque gaz. Quels gaz s'écartent le plus du modèle ? Proposer une explication à
l'échelle moléculaire.
\item Un extincteur de \qty{3,0}{L} contient \qty{2,0}{kg} de dioxyde de carbone à \qty{20}{\degreeCelsius}. Quelle pression le
modèle du gaz parfait prévoit-il ?
\item La pression réelle n'est que de \qty{57}{bar}. Expliquer.
\end{questions}
\end{exercice}
\end{blocLivre}
\rubrique{Objectif bac}
\begin{blocLivre}{ex-18-12}
\begin{exercice}[3]{Le coussin gonflable de sécurité}
\begin{docu}{1}{Un gonflage en 30 millisecondes}
Lors d'un choc, un capteur déclenche la décomposition d'un solide, l'azoture de sodium \ce{NaN3} ; le diazote libéré
gonfle le coussin en une trentaine de millisecondes :
\[\ce{2 NaN3(s) -> 2 Na(s) + 3 N2(g)}.\]
Toxique, l'azoture a depuis été remplacé par d'autres générateurs de gaz.
\end{docu}
\begin{docu}{2}{Le coussin gonflé}
Volume : \qty{60}{L} ; pression du gaz : \qty{1,30}{bar} ; température du gaz : \qty{50}{\degreeCelsius}.
\end{docu}
\begin{donnees}
$R=\qty{8,314}{J.mol^{-1}.K^{-1}}$ ; $M(\ce{NaN3})=\qty{65,0}{g.mol^{-1}}$.
\end{donnees}
\begin{questions}
\item Calculer la quantité de diazote contenue dans le coussin gonflé.
\item En déduire la quantité d'azoture de sodium qui doit réagir.
\item \initiative Quelle masse d'azoture de sodium faut-il prévoir ?
\item Le gaz se refroidit ensuite à \qty{20}{\degreeCelsius}. Quelle serait sa pression si le coussin restait fermé, de volume
constant ? En réalité, des évents le dégonflent : pourquoi est-ce indispensable ?
\end{questions}
\end{exercice}
\end{blocLivre}
\begin{blocLivre}{ex-18-13}
\begin{exercice}[3]{Le ballon-sonde éclate}
\begin{docu}{1}{Un ballon qui gonfle en montant}
Un ballon-sonde en latex est gonflé au sol avec \qty{3,0}{m^3} d'hélium, sous \qty{1013}{hPa} et à \qty{15}{\degreeCelsius}.
L'enveloppe, très souple, laisse le gaz à la pression extérieure ; elle éclate lorsque son diamètre atteint
\qty{8,0}{m}.
\end{docu}
\begin{docu}{2}{L'atmosphère standard}
\begin{center}
\begin{tabular}{lcccc}
\toprule
Altitude (km) & 20 & 25 & 30 & 35 \\
Pression (hPa) & 54,7 & 25,5 & 11,97 & 5,75 \\
Température (K) & 216,7 & 221,6 & 226,5 & 236,5 \\
\bottomrule
\end{tabular}
\end{center}
\end{docu}
\begin{donnees}
$R=\qty{8,314}{J.mol^{-1}.K^{-1}}$ ; volume d'une sphère de diamètre $d$ : $\dfrac{\pi\,d^3}{6}$. On admet que l'hélium reste
à la température de l'air extérieur.
\end{donnees}
\begin{questions}
\item Calculer la quantité d'hélium contenue dans le ballon.
\item Calculer le volume du ballon à \qty{30}{km} d'altitude, puis son diamètre. Le comparer au diamètre au sol.
\item \initiative Entre quelles altitudes le ballon éclate-t-il ?
\item Pourquoi le modèle du gaz parfait est-il ici particulièrement bien adapté ?
\end{questions}
\end{exercice}
\end{blocLivre}
\begin{blocLivre}{ex-18-14}
\begin{exercice}[3]{Respirer au sommet de l'Everest}
\begin{docu}{1}{Le toit du monde}
Au sommet de l'Everest, à \qty{8849}{m}, la pression atmosphérique n'est que de \qty{337}{hPa} et la température vaut
souvent \qty{-25}{\degreeCelsius} pendant la saison des ascensions. Au niveau de la mer, on prendra \qty{1013}{hPa} et
\qty{15}{\degreeCelsius}. L'air contient 21~\% de dioxygène en quantité de matière, quelle que soit l'altitude.
\end{docu}
\begin{docu}{2}{Dans les poumons}
L'air inspiré est réchauffé à \qty{37}{\degreeCelsius} dans les poumons, à la pression ambiante.
\end{docu}
\begin{donnees}
$R=\qty{8,314}{J.mol^{-1}.K^{-1}}$ ; $M(\text{air})=\qty{29,0}{g.mol^{-1}}$.
\end{donnees}
\begin{questions}
\item Calculer la masse volumique de l'air au niveau de la mer et au sommet.
\item Calculer la quantité de dioxygène contenue dans un litre d'air extérieur, au niveau de la mer puis au sommet.
\item \initiative Un alpiniste remplit ses poumons du même volume au niveau de la mer et au sommet. Montrer que la quantité
de dioxygène inspirée au sommet n'est que le tiers environ de celle inspirée au niveau de la mer.
\item Pourquoi la plupart des alpinistes utilisent-ils des bouteilles de dioxygène au-delà de \qty{8000}{m} ?
\end{questions}
\end{exercice}
\end{blocLivre}
\begin{blocLivre}{ex-18-15}
\begin{exercice}[3]{Pourquoi un moteur Diesel n'a pas de bougie}
\begin{docu}{1}{L'allumage par compression}
Dans un moteur Diesel, le piston comprime fortement l'air admis dans le cylindre ; le gazole, injecté à la fin de la
compression, s'enflamme spontanément au contact de l'air chaud. Le gazole s'auto-enflamme au-dessus de
\qty{250}{\degreeCelsius} environ.
\end{docu}
\schema{18-ex15}
\begin{docu}{2}{Mesures}
Au début de la compression : \qty{500}{cm^3} d'air, \qty{1,0}{bar}, \qty{30}{\degreeCelsius}. À la fin : volume divisé par 16
(rapport volumétrique), pression \qty{45}{bar}.
\end{docu}
\begin{donnees}
$R=\qty{8,314}{J.mol^{-1}.K^{-1}}$ ; $M(\text{air})=\qty{29,0}{g.mol^{-1}}$ ; on néglige les fuites.
\end{donnees}
\begin{questions}
\item Calculer la quantité et la masse d'air enfermées dans le cylindre.
\item Quelle pression obtiendrait-on en fin de compression si la température restait constante ?
\item \initiative Calculer la température de l'air en fin de compression. Pourquoi le gazole s'enflamme-t-il sans
étincelle ?
\item Dans un moteur à essence, le rapport volumétrique ne dépasse guère 10 et une bougie déclenche la combustion.
Pourquoi ne comprime-t-on pas davantage le mélange d'air et d'essence ?
\end{questions}
\end{exercice}
\end{blocLivre}
